\documentclass[12pt]{article}

\usepackage{amsmath}
\usepackage{amssymb}
\usepackage{physics}
\usepackage{hyperref}

\title{Affine Subspaces and Coordinates}
\author{}
\date{}

\begin{document}

\maketitle

\section{Affine Subspaces}

\begin{definition}[Affine subspace]
Let \(x\in \mathcal A\), and let \(U\) be a vector subspace of \(V\). The affine
subspace of \(\mathcal A\) determined by \(x\) and \(U\) is
\[
S(x,U)
=
\left\{
\mathbf{u}(x)
\;\middle|\;
\mathbf{u}\in U
\right\}.
\]
Its dimension is defined by
\[
\dim S(x,U)=\dim U.
\]
\end{definition}

\begin{proposition}
An \(n\)-dimensional affine space \(\mathcal A\) has only one \(n\)-dimensional
affine subspace, namely \(\mathcal A\) itself.
\end{proposition}

\begin{proof}
Suppose that
\[
\dim U=n.
\]
Let
\[
\mathbf{u}_1,\ldots,\mathbf{u}_n\in U
\]
be a basis of \(U\). Thus,
\[
\lambda_1\mathbf{u}_1+\cdots+\lambda_n\mathbf{u}_n=\mathbf{0}
\quad\Longleftrightarrow\quad
\lambda_i=0
\quad\text{for every }i=1,\ldots,n.
\]
Since \(U\) and \(V\) both have dimension \(n\), we have
\[
\operatorname{span}
\left\{
\mathbf{u}_1,\ldots,\mathbf{u}_n
\right\}
=
U
=
V.
\]

Fix \(x\in \mathcal A\). We know that for every \(y\in \mathcal A\), there exists a unique
vector \(\mathbf{u}\in V\) such that
\[
\mathbf{u}(x)=y.
\]
Since \(U=V\), we have \(\mathbf{u}\in U\). Therefore,
\[
y\in S(x,U).
\]
Hence,
\[
S(x,U)=\mathcal A.
\]
\end{proof}

\begin{proposition}
Let \(S(x,U)\) be an affine subspace determined by the point \(x\) and
the linear subspace \(U\). Then
\[
U
=
\left\{
\overrightarrow{yz}
\;\middle|\;
y,z\in S(x,U)
\right\}.
\]
\end{proposition}

\begin{proof}
Let us first consider \(\mathbf{u}\in U\), and let
\[
y=\mathbf{u}(x).
\]
By the definition of \(S(x,U)\), we have \(y\in S(x,U)\). Moreover,
\[
\mathbf{u}=\overrightarrow{xy}.
\]
Since \(x,y\in S(x,U)\), it follows that
\[
U
\subseteq
\left\{
\overrightarrow{yz}
\;\middle|\;
y,z\in S(x,U)
\right\}.
\]

We now prove the opposite inclusion. Let \(y,z\in S(x,U)\). By the
definition of \(S(x,U)\), there exist \(\mathbf{u},\mathbf{v}\in U\)
such that
\[
y=\mathbf{u}(x)
\qquad\text{and}\qquad
z=\mathbf{v}(x).
\]
Then
\[
(\mathbf{v}-\mathbf{u})(y)
=
(\mathbf{v}-\mathbf{u})\bigl(\mathbf{u}(x)\bigr)
=
\bigl((\mathbf{v}-\mathbf{u})+\mathbf{u}\bigr)(x)
=
\mathbf{v}(x)
=
z.
\]
Therefore,
\[
\overrightarrow{yz}
=
\mathbf{v}-\mathbf{u}.
\]
Since \(\mathbf{u},\mathbf{v}\in U\) and \(U\) is a vector subspace,
\[
\mathbf{v}-\mathbf{u}\in U.
\]
It follows that
\[
\left\{
\overrightarrow{yz}
\;\middle|\;
y,z\in S(x,U)
\right\}
\subseteq U.
\]
Hence,
\[
U
=
\left\{
\overrightarrow{yz}
\;\middle|\;
y,z\in S(x,U)
\right\}.
\]
\end{proof}

\begin{proposition}
Let \(x,y\in \mathcal A\), and let \(U,W\) be linear subspaces of \(V\). Then
\[
S(x,U)=S(y,W)
\]
if and only if
\[
U=W
\qquad\text{and}\qquad
S(x,U)\cap S(y,W)\neq\varnothing.
\]
\end{proposition}

\begin{proof}
Let us first assume that
\[
S(x,U)=S(y,W).
\]
It is then obvious that
\[
S(x,U)\cap S(y,W)\neq\varnothing.
\]
We need to prove that \(U=W\).

Let \(\mathbf{u}\in U\). There exists \(z\in S(x,U)\) such that
\[
z=\mathbf{u}(x).
\]
Since \(S(x,U)=S(y,W)\), we also have \(z\in S(y,W)\). Therefore, there
exists \(\mathbf{v}\in W\) such that
\[
z=\mathbf{v}(y).
\]
It follows that
\[
\mathbf{u}(x)=\mathbf{v}(y),
\]
and therefore
\[
x=(\mathbf{v}-\mathbf{u})(y).
\]

Since
\[
x\in S(x,U)=S(y,W),
\]
the vector mapping \(y\) onto \(x\) belongs to \(W\). Hence, there exists
\(\mathbf{w}\in W\) such that
\[
\mathbf{v}-\mathbf{u}=\mathbf{w}.
\]
Therefore,
\[
\mathbf{u}=-\mathbf{w}+\mathbf{v}.
\]
Since \(\mathbf{v},\mathbf{w}\in W\), it follows that
\[
\mathbf{u}\in W.
\]
Thus,
\[
U\subseteq W.
\]
Similarly, we can prove that
\[
W\subseteq U.
\]
Therefore,
\[
U=W.
\]

Conversely, let us assume that
\[
U=W
\qquad\text{and}\qquad
S(x,U)\cap S(y,W)\neq\varnothing.
\]
Let
\[
z\in S(x,U)\cap S(y,W).
\]
Then there exist \(\mathbf{u}\in U\) and \(\mathbf{v}\in W\) such that
\[
z=\mathbf{u}(x)=\mathbf{v}(y).
\]

We also know that for every \(k\in S(x,U)\), there exists
\(\mathbf{q}\in U\) such that
\[
k=\mathbf{q}(x).
\]
Since
\[
\mathbf{u}(x)=\mathbf{v}(y),
\]
we have
\[
x=(\mathbf{v}-\mathbf{u})(y).
\]
Therefore,
\[
k
=
\mathbf{q}(x)
=
\mathbf{q}\bigl((\mathbf{v}-\mathbf{u})(y)\bigr)
=
(\mathbf{q}+\mathbf{v}-\mathbf{u})(y).
\]
Since \(U=W\), we have
\[
\mathbf{q}+\mathbf{v}-\mathbf{u}\in W.
\]
It follows that
\[
k\in S(y,W).
\]
Therefore,
\[
S(x,U)\subseteq S(y,W).
\]
Similarly, we can prove that
\[
S(y,W)\subseteq S(x,U).
\]
Hence,
\[
S(x,U)=S(y,W).
\]
\end{proof}
\begin{proposition}
Let \(S(x,U)\) be a \(r\)-dimensional affine subspace of \(\mathcal A\). Then,
with the action inherited from \(\mathcal A\), \(S(x,U)\) is itself a
\(r\)-dimensional affine space over the vector space \(U\).
\end{proposition}

\begin{proof}
Since \(U\) is a vector subspace of \(V\), the action of \(U\) on
\(S(x,U)\) is well-defined. Indeed, let \(\mathbf{u}\in U\) and
\(y\in S(x,U)\). Then there exists \(\mathbf{p}\in U\) such that
\[
y=\mathbf{p}(x).
\]
Therefore,
\[
\mathbf{u}(y)
=
\mathbf{u}\bigl(\mathbf{p}(x)\bigr)
=
(\mathbf{u}+\mathbf{p})(x).
\]
Since
\[
\mathbf{u}+\mathbf{p}\in U,
\]
it follows that
\[
\mathbf{u}(y)\in S(x,U).
\]

We know that \(U\) is a vector space. Hence, it has the null vector
\[
\mathbf{0}\in U,
\]
and, for every \(y\in S(x,U)\),
\[
\mathbf{0}(y)=y.
\]

We now need to prove that, for every \(z,y\in S(x,U)\), there exists a
unique \(\mathbf{u}\in U\) such that
\[
\mathbf{u}(y)=z.
\]
We know that there exist unique vectors \(\mathbf{q},\mathbf{p}\in U\)
such that
\[
z=\mathbf{q}(x)
\qquad\text{and}\qquad
y=\mathbf{p}(x).
\]
Therefore,
\[
x=(-\mathbf{p})(y),
\]
and hence
\[
z
=
\mathbf{q}(x)
=
\mathbf{q}\bigl((-\mathbf{p})(y)\bigr)
=
(\mathbf{q}-\mathbf{p})(y),
\]
because of the inherited composition.

Therefore, there exists a unique \(\mathbf{u}\in U\) such that
\[
z=\mathbf{u}(y),
\]
where
\[
\mathbf{u}=\mathbf{q}-\mathbf{p}.
\]

Finally, the composition property is inherited from the action of \(V\)
on \(\mathcal A\). For every \(\mathbf{u},\mathbf{v}\in U\) and
\(y\in S(x,U)\),
\[
(\mathbf{u}\circ\mathbf{v})(y)
=
\mathbf{u}\bigl(\mathbf{v}(y)\bigr)
=
(\mathbf{u}+\mathbf{v})(y).
\]

Therefore, \(S(x,U)\) is an affine space over \(U\). Since
\[
\dim U=r,
\]
the affine space \(S(x,U)\) has dimension \(r\).
\end{proof}

\begin{proposition}
Let \(S(x,U)\) and \(S(y,W)\) be affine subspaces of \(\mathcal A\). If
\[
z\in S(x,U)\cap S(y,W),
\]
then
\[
S(x,U)\cap S(y,W)
=
S(z,U\cap W).
\]
\end{proposition}

\begin{proof}
Since
\[
z\in S(x,U)\cap S(y,W),
\]
there exist \(\mathbf{u}\in U\) and \(\mathbf{v}\in W\) such that
\[
z=\mathbf{u}(x)=\mathbf{v}(y).
\]

We first consider an arbitrary point
\[
h\in S(x,U)\cap S(y,W).
\]
Then there exist \(\mathbf{q}\in U\) and \(\mathbf{p}\in W\) such that
\[
h=\mathbf{q}(x)=\mathbf{p}(y).
\]

Since
\[
z=\mathbf{u}(x),
\]
we have
\[
x=(-\mathbf{u})(z),
\]
and therefore
\[
h
=
\mathbf{q}(x)
=
\mathbf{q}\bigl((-\mathbf{u})(z)\bigr)
=
(\mathbf{q}-\mathbf{u})(z).
\]
Similarly, since
\[
z=\mathbf{v}(y),
\]
we have
\[
h
=
\mathbf{p}(y)
=
\mathbf{p}\bigl((-\mathbf{v})(z)\bigr)
=
(\mathbf{p}-\mathbf{v})(z).
\]

By the uniqueness of the vector mapping \(z\) onto \(h\),
\[
\mathbf{q}-\mathbf{u}
=
\mathbf{p}-\mathbf{v}
=
\mathbf{m}.
\]
Since
\[
\mathbf{q}-\mathbf{u}\in U
\qquad\text{and}\qquad
\mathbf{p}-\mathbf{v}\in W,
\]
it follows that
\[
\mathbf{m}\in U\cap W.
\]
Hence,
\[
h=\mathbf{m}(z),
\]
and therefore
\[
S(x,U)\cap S(y,W)
\subseteq
S(z,U\cap W).
\]

We now take a point
\[
k\in S(z,U\cap W).
\]
Then there exists \(\mathbf{q}\in U\cap W\) such that
\[
k=\mathbf{q}(z).
\]
Since
\[
z=\mathbf{u}(x)=\mathbf{v}(y),
\]
we have
\[
k
=
\mathbf{q}(z)
=
\mathbf{q}\bigl(\mathbf{u}(x)\bigr)
=
(\mathbf{q}+\mathbf{u})(x),
\]
and
\[
k
=
\mathbf{q}(z)
=
\mathbf{q}\bigl(\mathbf{v}(y)\bigr)
=
(\mathbf{q}+\mathbf{v})(y).
\]
Since
\[
\mathbf{q}+\mathbf{u}\in U
\qquad\text{and}\qquad
\mathbf{q}+\mathbf{v}\in W,
\]
we have
\[
k\in S(x,U)\cap S(y,W).
\]
Therefore,
\[
S(z,U\cap W)
\subseteq
S(x,U)\cap S(y,W).
\]

Hence,
\[
S(x,U)\cap S(y,W)
=
S(z,U\cap W).
\]
\end{proof}

\begin{definition}[Coordinate system]
Let \(S(x,U)\) be an affine space of dimension \(r\). Let
\[
\mathbf{u}_1,\mathbf{u}_2,\ldots,\mathbf{u}_r
\]
be a basis for \(U\), and let \(x_0\) be a point in \(S(x,U)\). The point \(x_0\) and the basis \(\mathbf{u}_1,\ldots,\mathbf{u}_r\) form a coordinate system for \(S(x,U)\).

We use the Einstein summation convention throughout: an index that occurs twice in a term is implicitly summed over its full range.

Let \(y\in S(x,U)\). There exist unique vectors
\(\mathbf{u},\mathbf{u}_0\in U\) such that
\[
y=\mathbf{u}(x)
\qquad\text{and}\qquad
x_0=\mathbf{u}_0(x).
\]
Hence
\[
x=(-\mathbf{u}_0)(x_0),
\]
and therefore
\[
y
=
\mathbf{u}\bigl((-\mathbf{u}_0)(x_0)\bigr)
=
(\mathbf{u}-\mathbf{u}_0)(x_0).
\]
Thus
\[
\overrightarrow{x_0y}=\mathbf{u}-\mathbf{u}_0.
\]
Writing
\[
\mathbf{u}-\mathbf{u}_0=y^i\mathbf{u}_i,
\]
the components \(y^1,\ldots,y^r\) are called the coordinates of \(y\) with respect to this coordinate system.
\end{definition}

\begin{proposition}
Given a coordinate system \(x_0,\mathbf{u}_i\) for an affine subspace
\(S(x,U)\), the mapping
\[
c:S(x,U)\longrightarrow\mathbb{R}^r,
\qquad
c(y)=\left(y^1,y^2,\ldots,y^r\right),
\]
is one-to-one and onto.
\end{proposition}

\begin{proof}
First, we note that \(c\) is the composition of two mappings:
\[
S(x,U)\xrightarrow{\,g\,}U\xrightarrow{\,h\,}\mathbb{R}^r.
\]

For any \(y\in S(x,U)\), let
\[
y=\mathbf{u}(x)
\qquad\text{and}\qquad
x_0=\mathbf{u}_0(x).
\]
We define
\[
g(y)=\mathbf{u}-\mathbf{u}_0.
\]
If
\[
\mathbf{v}=v^i\mathbf{u}_i,
\]
we define
\[
h(\mathbf{v})=\left(v^1,v^2,\ldots,v^r\right).
\]
Therefore,
\[
c(y)=(h\circ g)(y).
\]

Let us first prove that \(c\) is one-to-one. Assume that
\[
(h\circ g)(z)=(h\circ g)(y).
\]
Let
\[
z=\mathbf{q}(x)
\qquad\text{and}\qquad
y=\mathbf{u}(x).
\]
Then
\[
g(z)=\mathbf{q}-\mathbf{u}_0
\qquad\text{and}\qquad
g(y)=\mathbf{u}-\mathbf{u}_0.
\]
Hence
\[
h(\mathbf{q}-\mathbf{u}_0)
=
h(\mathbf{u}-\mathbf{u}_0).
\]
The components of a vector with respect to a basis are unique, and therefore
\[
\mathbf{q}-\mathbf{u}_0
=
\mathbf{u}-\mathbf{u}_0.
\]
Thus
\[
\mathbf{q}=\mathbf{u},
\]
and consequently
\[
z=y.
\]
Hence \(c\) is one-to-one.

Let us now prove that \(c\) is onto. We first prove that
\[
h:U\longrightarrow\mathbb{R}^r
\]
is onto. Indeed, for any
\[
\mathbf{k}=\left(k^1,k^2,\ldots,k^r\right)\in\mathbb{R}^r,
\]
there exists a vector \(\mathbf{p}\in U\) given by
\[
\mathbf{p}=k^i\mathbf{u}_i,
\]
such that
\[
h(\mathbf{p})=\mathbf{k}.
\]

We now prove that
\[
g:S(x,U)\longrightarrow U
\]
is onto. Indeed, for any \(\mathbf{p}\in U\), set
\[
\mathbf{u}=\mathbf{p}+\mathbf{u}_0.
\]
Then \(\mathbf{u}\in U\), and
\[
y=(\mathbf{u}-\mathbf{u}_0)(x_0)\in S(x,U).
\]
Moreover,
\[
g(y)=\mathbf{u}-\mathbf{u}_0=\mathbf{p}.
\]
Therefore \(g\) is onto.

Since both \(g\) and \(h\) are onto, their composition
\[
c=h\circ g
\]
is onto.
\end{proof}

\end{document}
